\chapter{产品化、迭代与实战}
\label{chap:edge-ai-product-lifecycle-projects}

模型第一次在板端运行只是生命周期的中点。量产系统还要面对新设备、新批次、光照变化、数据漂移、版本组合、远程升级、隐私和现场回退。本章把前述数据、模型、artifact 和 runtime 组织为可持续演进的产品系统，并给出从入门到精通的综合项目路线。

\section{把模型当作版本化产品组件}

模型更新至少同时涉及：

\begin{itemize}
  \item checkpoint、portable graph、target artifact 和 Manifest；
  \item label/schema、预处理、参考解码、阈值和产品策略；
  \item converter/runtime/driver/firmware 与硬件 profile；
  \item 校准、测试集、gold inputs 和性能基线；
  \item 已持久化的 Embedding、tracker 状态或派生记录。
\end{itemize}

只替换一个模型文件可能造成类别顺序、输入颜色、量化 scale 或数据库向量版本不兼容。发布包应声明兼容矩阵，加载时拒绝未知组合。

\section{模型发布状态机}

推荐把模型从训练到产品的状态显式化：

\begin{center}
\texttt{experiment $\rightarrow$ evaluated $\rightarrow$ converted}
\\
\texttt{$\rightarrow$ board-verified $\rightarrow$ staged $\rightarrow$ active $\rightarrow$ retired}
\end{center}

状态转换需要证据和责任人：

\begin{description}
  \item[evaluated] 独立数据、困难切片和预注册质量门通过；
  \item[converted] portable graph、量化和 artifact metadata 通过；
  \item[board-verified] 真实目标硬件完成一致性、性能和稳定验证；
  \item[staged] 已部署到受控设备，具备监控和自动/人工回退；
  \item[active] 成为批准默认版本；
  \item[retired] 不再新部署，但保留审计与必要回滚信息。
\end{description}

训练成功、转换成功和产品晋级是三个不同动作。

\section{模型 CI/CD 与硬件在环流水线}

模型流水线不应把所有工作塞进一次超长 CI。快速、确定的检查在每次变更运行；GPU 训练、厂商转换、板端长稳和灰度发布由版本化输入触发，并保留人工批准边界。

\begin{table}[H]
  \centering
  \caption{边缘 AI 发布流水线的分层 Gate}
  \begin{tabular}{p{26mm}p{55mm}p{51mm}}
    \hline
    层级 & 自动化检查 & 失败处理 \\
    \hline
    Source/Data & schema、manifest、重复/group 泄漏、许可证与格式 & 阻止产生新数据版本 \\
    Model unit & Target/Loss、valid 零梯度、增强往返、单 batch 与小样本过拟合 & 修复实现，不启动正式训练 \\
    Training & 冻结配置、唯一 run、完成 step、指标与 checkpoint 摘要 & 标记 run 失败或不完整 \\
    Export & graph check、静态 Shape、Manifest、L0 parity & 不进入 converter \\
    Conversion & profile/version、日志、artifact metadata、L2/L3 & 不进入板端默认路径 \\
    Board/HIL & gold tensor、ABI、P99/内存/功耗、故障注入和长稳 & 保留 incumbent \\
    Release & 签名、兼容矩阵、灰度 guardrail 与 rollback rehearsal & 停止或自动回退 \\
    \hline
  \end{tabular}
\end{table}

一次流水线运行需要把以下身份绑定在 promotion record 中：source revision、数据/schema/split、run/checkpoint、portable graph、converter/profile、target artifact、runtime/driver/firmware、评估报告和设备结果。Artifact registry 或对象存储只保存文件还不够，必须能够从任一产物反查完整依赖和批准状态。

GPU 与板端资源通常稀缺，队列应支持互斥、超时、取消、重试和结果去重。重试不能覆盖原 run；同一输入摘要已成功产生的确定性 artifact 可以复用，但训练或硬件结果必须保留实际执行身份。板卡农场还要记录板型/序列号、电源、温度、频率、SDK/runtime 和复位状态，避免把环境差异误判为模型回归。

推荐把发布权限与训练权限分开。训练系统可以产生 candidate，但只有通过独立评估、板端 Gate 和审批的不可变包才能进入 staged/active。流水线任何未执行层级都标为 unavailable，不能用绿色的上游任务替代缺失的目标板证据。

\section{安全更新与回滚}

模型更新沿用系统 OTA 的完整性、原子性和防回滚机制，参见第~\ref{chap:system-build-update}~章。模型特有检查包括：

\begin{enumerate}
  \item 下载前检查硬件、runtime、内存和存储适用性；
  \item 验证签名、摘要、Manifest schema 和 ABI；
  \item 在候选路径加载并运行固定 gold input；
  \item 灰度启用，监控错误、latency、资源和业务 guardrail；
  \item 健康确认后切换 canonical 版本；
  \item 失败时原子回退模型、decoder、阈值和相关配置；
  \item 保留失败原因和设备环境，避免无限重复升级。
\end{enumerate}

如果模型 ABI 改变，可并行保留旧/新 adapter。不能用启发式根据输出数值猜版本。

\section{现场监控不等于上传原始图像}

默认监控应优先使用不含原始敏感内容的统计：

\begin{itemize}
  \item 输入尺寸、亮度/颜色分布、模糊、遮挡和域质量；
  \item 类别/unknown/empty 分布、score 和 margin；
  \item 每阶段 latency、队列深度、drop、timeout 和 reset；
  \item 模型/runtime/hardware/calibration generation；
  \item 资源、温度、频率和内存；
  \item 经授权抽样的错误证据或脱敏特征。
\end{itemize}

监控指标必须有基线、窗口和动作。类别比例变化可能来自真实业务变化，也可能是相机偏移、照明损坏、模型漂移或数据上报故障，不能自动归因。

原始图像、音视频和人脸/身份 Embedding 具有更高隐私风险。采集、上传、保留和访问应最小化并进入安全生命周期；崩溃转储默认不包含完整 tensor 或媒体。

\section{漂移、反馈与新数据}

漂移可以分为：

\begin{description}
  \item[输入漂移] 相机、光照、背景、对象或 ISP 分布变化；
  \item[标签先验漂移] 各状态出现频率变化；
  \item[概念漂移] 相同视觉证据对应的业务含义变化；
  \item[系统漂移] runtime、驱动、预处理或硬件版本改变；
  \item[观测漂移] 日志/采样策略改变导致监控分布变化。
\end{description}

检测漂移不等于证明质量下降。需要把漂移样本进入人工审核或独立金标回归，再判断采集、标注、阈值、模型还是系统路径需要更新。

一个健康反馈循环为：

\begin{enumerate}
  \item 现场指标或人工反馈发现高价值切片；
  \item 按隐私和授权采集代表样本；
  \item 去重、分组并进入审核队列；
  \item 更新 reviewed 数据版本；
  \item 在冻结 incumbent 和评估协议下训练 candidate；
  \item 完成导出、量化、板端和灰度 Gate；
  \item 仅在证据通过后晋级。
\end{enumerate}

\section{主动学习与人工纠正保护}

成熟系统用不确定性、教师分歧、稀有性、域新颖度和业务价值选择审核样本。人工纠正是最高优先级事实，后续预标注不得覆盖。自动接受阈值按类别校准，并持续抽检；错误率或域分布超限时撤销自动放行。

每轮主动学习报告人工小时、接受/修改/漏标率、每任务新增有效字段和固定回归集增益。当新增标注不再提升时，应检查标签歧义、域覆盖和模型容量，而不是继续积累相似图片。

\section{多模型、级联与时间状态}

高级产品常见拓扑：

\begin{itemize}
  \item 检测器产生 ROI，分类/关键点模型二次分析；
  \item 多个并行模型融合检测、分割和质量；
  \item 跟踪器和时序状态机把单帧证据转换为稳定事件；
  \item 多摄像头或视觉与深度/雷达融合；
  \item 高精度 Teacher 在服务器预标注，轻量 Student 在边缘运行。
\end{itemize}

级联必须保存每次 crop、resize、padding 和逆映射，并定义最大 ROI、批处理、调度和旧结果失效。不同对象、帧、stream 或模型版本的输出不能误关联。

时间策略需要 timestamp 单位、乱序、丢帧、EMA/投票、迟滞、debounce、冷却、reset 和 stream 隔离。业务事件属于产品策略，训练指标之外还要评估事件级 Precision/Recall、响应时间和持续误报。

\section{安全、鲁棒与越界输入}

视觉系统需要防御：

\begin{itemize}
  \item 损坏或超大媒体导致 parser/decoder 资源耗尽；
  \item 非法 tensor metadata、Shape、stride 和地址；
  \item 模型/配置篡改、降级和未授权替换；
  \item 对抗贴纸、强光、遮挡和故意构造背景；
  \item 域外输入被高置信接受；
  \item 日志、抓图、Embedding 和调试接口泄露敏感数据。
\end{itemize}

安全不是要求模型在所有攻击下正确，而是定义攻击面、检测、拒绝、降级和安全状态。高影响动作不能仅由一个未经质量门控的视觉输出直接触发；可结合多传感器、物理联锁或人工确认。

模型卡和数据卡应说明用途、限制和不建议用途。AI 风险治理可参考 NIST AI RMF 的治理、映射、测量和管理思路\cite{nist-ai-rmf}，具体产品仍需结合安全、法规和行业标准。

\section{成本、容量与长期维护}

产品设计不仅优化每帧 latency，还要计算：

\begin{itemize}
  \item 数据采集、存储、标注和审核成本；
  \item GPU 训练、实验失败和 artifact 保留成本；
  \item 多芯片 converter/runtime 维护矩阵；
  \item 设备端模型存储、下载流量、升级窗口和回滚空间；
  \item 现场监控、人工复核和隐私合规成本；
  \item 供应商 SDK 生命周期、芯片替代和模型重新认证。
\end{itemize}

新抽象、backend 或模型 profile 必须有当前消费者和可测收益。每个实验开关长期保留，会扩大测试组合并增加现场不确定性。

\section{综合项目一：入门分类器}

\textbf{目标：} 在固定 ROI 上完成四分类与 unknown 拒绝。

\textbf{必须交付：}

\begin{itemize}
  \item 三个以上 session 的 group-safe 数据版本；
  \item 迁移学习 baseline、小样本过拟合和混淆矩阵；
  \item portable graph、Manifest、reference inference 和 golden inputs；
  \item CPU 或开发板推理 latency、内存和错误路径；
  \item 一页模型卡，说明不支持场景。
\end{itemize}

\textbf{能力门：} 能解释 softmax/sigmoid、过拟合、阈值和框架到 portable 对齐。

\section{综合项目二：端侧目标检测}

\textbf{目标：} 在一个画面中检测物料和方向错误，接入实时视频。

\textbf{必须交付：}

\begin{itemize}
  \item box schema、标注一致性、空场景和强负样本；
  \item 同一 harness 下至少两个 model scale 的比较；
  \item 每类 AP/Recall、小目标/遮挡/设备切片和 failure gallery；
  \item PTQ、L0/L2/L3 证据和目标 artifact；
  \item Camera-to-runtime pipeline、gold tensor、frame identity 和 P99；
  \item timeout、队列满、摄像头断流和加速器 reset 注入。
\end{itemize}

\textbf{能力门：} 能定位预处理、导出、量化、后处理和运行时的首个分歧。

\section{综合项目三：分割/关键点与多阶段系统}

\textbf{目标：} 对细小缺陷或装配几何进行分割/关键点分析，并与检测 ROI 级联。

\textbf{必须交付：}

\begin{itemize}
  \item mask/关键点 valid、不可见和坐标变换测试；
  \item foreground/boundary 或关键点误差切片；
  \item 级联 ROI、最大候选、调度和关联规则；
  \item 多模型内存、DDR、latency、热稳态和降级 profile；
  \item 产品状态机、旧结果失效和事件级评估。
\end{itemize}

\textbf{能力门：} 能把模型语义、backend 机制和产品策略分层，并证明跨阶段对象/坐标没有混淆。

\section{综合项目四：可持续模型产品}

\textbf{目标：} 建立一次完整的领域数据反馈、候选训练、灰度发布和回滚演练。

\textbf{必须交付：}

\begin{itemize}
  \item 现场漂移指标和受控采集/审核流程；
  \item Teacher 预标注、人工纠正保护和主动学习队列；
  \item 冻结 incumbent、预注册候选规则和多切片评估；
  \item 多 backend Artifact Manifest、签名、兼容矩阵和 SBOM 关联；
  \item staged rollout、健康门、自动/人工 rollback；
  \item 发布后监控、残余风险和下一轮数据计划。
\end{itemize}

\textbf{能力门：} 能把一次模型实验转化为可审计、可回退、可长期维护的产品变更。

\section{建议的学习节奏}

对具备嵌入式软件基础的读者，可以按以下节奏安排：

\begin{table}[H]
  \centering
  \caption{边缘视觉学习节奏示例}
  \begin{tabular}{p{22mm}p{43mm}p{68mm}}
    \hline
    周期 & 主题 & 阶段输出 \\
    \hline
    第 1--2 周 & 分类闭环与张量基础 & 数据 v001、baseline、portable inference \\
    第 3--5 周 & 数据 schema、检测与评估 & group split、detector、困难切片 \\
    第 6--8 周 & 导出、PTQ/QAT 与 manifest & L0--L3、模型包、拒绝/晋级记录 \\
    第 9--11 周 & Camera-to-runtime 与性能 & gold tensor、ABI tests、P99/内存/故障 \\
    第 12--16 周 & 多阶段、更新与主动学习 & 综合项目、灰度与回滚演练 \\
    持续阶段 & 新设备/域/芯片迁移 & profile 适配、现场反馈和长期证据 \\
    \hline
  \end{tabular}
\end{table}

时间只用于规划，阶段门仍以证据为准。已有深度学习背景者可以缩短前两阶段，但不能跳过数据切分、Manifest、量化对齐和故障路径。

\section{从熟练到精通的判断标准}

熟练工程师能够训练和部署一个已知模型；精通更体现在面对未知系统时能够：

\begin{itemize}
  \item 从业务失败语义选择最小充分视觉证据；
  \item 发现数据泄漏、标签歧义和不可观测问题；
  \item 冻结实验合同并区分真实收益与训练方差；
  \item 设计稳定 Model ABI、Manifest、参考解码和 backend profile；
  \item 用 L0--L4 找到第一个数值或系统分歧；
  \item 在质量、latency、内存、功耗、热和维护成本间做可审计权衡；
  \item 保持 backend、模型语义和产品策略分层；
  \item 设计现场监控、数据反馈、灰度、回滚和隐私边界；
  \item 对没有实际执行的 GPU、converter、simulator 或板端验证明确说“尚未验证”。
\end{itemize}

终点不是掌握某个框架或 NPU，而是把视觉证据转化为可验证、可交付、可演进的产品能力。
